% ============================================================
%  A4 - Serienbrief: ein Brief pro Zeile der CSV-Datei
%  Kompilieren: pdflatex serienbrief.tex
%  Benötigt im selben Ordner: fhbrief.cls, fh-logo.jpg, empfaenger.csv
% ============================================================
\documentclass[11pt]{fhbrief}
\usepackage{csvsimple}       % CSV-Dateien zeilenweise verarbeiten

\betreff{Einladung zum \LaTeX-Workshop}

\begin{document}

% \csvreader[Optionen]{Datei}{Spaltenzuordnung}{Inhalt pro Zeile}
%
% head to column names: Die erste Zeile der CSV-Datei liefert die
% Namen. Aus der Spalte "name" wird der Befehl \name, aus "ort"
% der Befehl \ort usw. Die Spaltenzuordnung {} bleibt daher leer.
%
% Der Inhalt in der letzten Klammer wird für jede Datenzeile einmal
% ausgeführt. Die Umgebung "brief" beginnt jedes Mal eine neue Seite.
\csvreader[head to column names]{empfaenger.csv}{}{%
  \empfaenger{\name\\\strasse\\\ort}%
  \begin{brief}{\anrede,}
    hiermit lade ich Sie zum \LaTeX-Workshop am Campus Hagenberg ein.

    Sie haben sich für das Thema \emph{\thema} interessiert.
    Dieses Thema ist Teil des Abschnitts \LaTeX{} Advanced.

    Ich freue mich auf Ihre Teilnahme.
  \end{brief}
}

\end{document}
